
\AddSymbolsH
\pagecolor{MidnightBlue!20!Gray}
\color{white}
\quad\quad \lettrine[,lraise=0.2,lhang=0.5]{\textbf{傍}}{晚}，天空格外明朗。

教授不安地走在河堤上。今天是夏至；空气炎热到可以把本应属于他的放松全部蒸发掉。

但奇怪的是，天空在下午五点时就已经变得漆黑。教授知道，自己离那个将要与他重逢的人愈来愈近了。

他拿出手里的两张卡纸\footnote{见第31-32页，以及第79-80页。}。他心中愧对自己对伊藤佐一做的一切——以及即将做的一切。

他焦灼地看向手表；已经逼近晚上九点了——可是他的焦灼似乎只持续了几分钟。

他突然停在原地。那个人就在前方。

是那位老者——让他茶饭不思、夜不能寐的老者。

他走近老者；可是连招呼都没打，时间已经开始逼近凌晨12点了。周围的人群载着孤独，在夜空下快速前进和后退着。

\twkkai{好久不见，老先生！}教授假装高兴地说着。

\twkkai{嗯。交给你的任务完成得怎么样了？}

老者推了推眼镜，极为平静地问着。

\twkkai{勉强完成了；}教授回答着。

老者转头看向教授。\twkkai{不必太认真的——我也只是希望自己死后的遗愿能够实现罢了；}老者说着。

教授更为不安地点了点头。他知道，老者的遗愿，便是他一生的诅咒。他决定在现在这个时刻把卡纸递给老者。

老者拿到卡纸后，像恼人的苏格拉底那般问起了第一个问题。

\twkkai{我活着的时候，最怕的不是死，是怕来不及……150年过去了，如今还有人认得出我吗？}

教授则像苏格拉底的学生柏拉图那样思考着，许久也不说话。

而他的身旁，突然多了一只橘猫，静静地尾随在老者身后。

教授没有留意到，而是像机器一样默背出了什么。猫猛地竖起耳朵，顺便也竖起了尾巴。

\twkkai{
您还记得您曾经定义过的密度吗？随着时代的发展，一切都已经变了；后来的数学家都非常感激您递给他们的接力棒，发明了代数数域的概念。一般的代数数域 \(F\) 的素理想集合 \(\mathcal{S}\) ，其用您名字命名的密度被定义为
\[
\delta_{F}(\mathcal{S})=\lim _{s \to 1^{+}} \frac{\sum_{\mathfrak{p} \in \mathcal{S}} \frac{1}{N \mathfrak{p}^{s}}}{\log \left(\frac{1}{s-1}\right)}.
\]}\index{leiyulun@类域论（Class field theory）!狄利克雷密度 $\delta_F(\mathcal{S}_{K/F})=1/[K:F]$}

老者的眼眶湿润了；他没想到，自己所做的一切，竟然给后世带去了这样的影响。而后世也竟然没有忘记他这个耄耋老人的功劳——他没有完全地离开这个世界。

教授像注意到证明条件一般瞥见了老者的泪珠。他不好做什么——怕影响老者的威严；只好继续开口说话——仿佛忘了自己的苦役。

但是，说着说着，他自己也哽咽起来——他感到一丝莫名其妙，甚至还为自己的哽咽冠上了自作多情之名。他依旧恭敬地默背着。

\twkkai{
现在……请允许我回到推广您在算术级数中素数定理证明里所用技巧的工作上来。群同态 \(\chi: I_{F}(\mathfrak{m}) / P_{F, \mathfrak{m}}^{+} \to \mathbb{C}^{\times}\)是被称为广义的以您的名字命名的特征，并且我们定义了韦伯\footnote{此处指数学家海因里希·马丁·韦伯（Heinrich Martin Weber, 1842–1913），区别于书中此前出现的作曲家卡尔·马利亚·冯·韦伯（Carl Maria von Weber）。} \(L\)-函数\[
L_{\mathfrak{m}}(s, \chi)=\sum_{\mathfrak{a} \in I(\mathfrak{m})} \frac{\chi(\mathfrak{a})}{N \mathfrak{a}^{s}}.
\]}\index{guangyitezhengbiao@广义狄利克雷特征}

老者推了推眼镜，用手指揩了揩眼泪，感叹着。

\twkkai{啊……我亲爱的\(L\)-函数啊——它长大啦……变得更强壮了，变得更秀丽了。看看它的下标\({\mathfrak{m}}\)吧……}

一旁四处张望的猫抬头看向老者。

\twkkai{喵……；}它小声地呜咽着——小声到教授和老者都根本无法听见。猫和老者都只顾沉浸在各自的感伤中。

教授仿佛才幡然醒悟。他的哽咽也越来越明显了；他终于意识到了什么——原来老者的威严背后，还有这样的一面。他彻底忘了他的苦役——他那测量月球半径的苦役。

天上的月亮眨了眨眼。

\twkkai{是的，是的——它长大了，而且依然具有欧拉乘积
\[
L_{\mathfrak{m}}(s, \chi)=\prod_{\mathfrak{p} \nmid \mathfrak{m}}(1-\chi(\mathfrak{p}) N \mathfrak{p}^{-s})^{-1}
\]
对 \(\operatorname{Re}(s)>1\) 成立。
}

教授顺着老者的话，温柔地陈述着。老者清了清喉咙，仿佛想说什么，但又不舍得让教授的言语消散。猫抓紧了脚步，跟了上来。风越来越大，直吹得猫发冷，让它连忙紧贴着老者的裤腿。

\twkkai{
后来，您启发了很多数学家。他们都断言，若对每个 \(R_{F, \mathfrak{m}}^{+}\) 的特征 \(\chi \neq \chi_{0}\) 都有 \(L_{\mathfrak{m}}(1, \chi) \neq 0\)，则
\[
\delta_{F}\left(\mathcal{S}_{\mathfrak{a}, \mathfrak{m}}\right)=\frac{1}{\# \mathcal{R}_{F, \mathfrak{m}}^{+}}.
\]

当然，这说明 \(\mathcal{S}_{\mathfrak{a}, \mathfrak{m}}\) 包含无穷多个素理想。所以……
}

教授终于顿了顿。老者竟然主动地把手搭在教授的肩膀上。

\twkkai{不哭，不哭——我的好孩子。}

教授哭得更厉害了。

猫见此，焦虑地在教授和老者两人的裤腿间画着“\(\infty\)”的轨迹。

\twkkai{所以……就是这样了。算术级数中素数定理的推广，将在我们证明了当 \(\chi \neq \chi_{0}\) 时 \(L_{\mathfrak{m}}(1, \chi) \neq 0\) 之后得到；}教授捂着脸说，痛哭着。

两个人，和一只猫，就那么停在河堤旁，挨在大理石围栏上，静静地看着河水东流。

它将载着什么而去，或许——连老者都不会知道。



